\documentclass[10pt,a4paper]{amsart}
\usepackage[margin=19mm]{geometry}
\usepackage{amsmath,amssymb,mathtools}
\usepackage[scr=rsfs,cal=euler]{mathalfa}
\usepackage{xfrac}
\usepackage[unicode,hidelinks,bookmarks=false]{hyperref}
\newcommand{\ie}{i.e.\ }
\newcommand{\cf}{cf.\ }
\newcommand{\resp}{respectively\ }
\newcommand{\Subsection}{\subsection}
\providecommand{\nobreakdash}{\nobreak\hbox{-}}
\setlength{\emergencystretch}{2em}
\multlinegap0pt
\usepackage{amssymb}
\usepackage{bm}
\newtheorem{theorem}{Theorem}
\title{Recursive Patterns in the Chocolate Game}
\author{Tomoro Okubo, Yuzuri Kashiwagi, Nobumitsu Niida}
\date{Source: arXiv:2602.20182v1 (CC BY 4.0)}
\begin{document}
\maketitle

\noindent\emph{Source and licence.} Recursive Patterns in the Chocolate Game. Version arXiv:2602.20182v1 (CC BY 4.0), distributed by arXiv under the Creative Commons Attribution 4.0 International licence, http://creativecommons.org/licenses/by/4.0/, as recorded in the arXiv licence metadata for this exact version and shown on the abstract page https://arxiv.org/abs/2602.20182v1. Full licence notice: \texttt{licenses/arxiv-2602.20182-CC-BY-4.0.txt}.\par\smallskip

\noindent\textit{Editorial scope.} Counted unit: the article's theorem theorem (source0.tex lines 602--629). The article's complete original proof of it is delivered with it (source0.tex lines 631--718, one \texttt{proof} environment of 88 lines). No mathematics is written, repaired or re-derived: the statement and the proof are the article's own lines, in the article's own notation, and no retained line is substituted or re-flowed.  No further source-local context is needed: every cross-reference of every retained line resolves inside the two delivered spans.  Nothing was omitted from any retained span, so the package carries no omission record.  No retained line contains a citation, so the wrapper carries no bibliography and no bibliography entry.  No retained line contains a figure environment, an image-inclusion command or drawing code, so no image is delivered and the package declares no asset dependency.  Editorial formatting only, in addition: an AMS \texttt{amsart} preamble that restates the article's own declarations and macros used by the retained lines, namely the environments \texttt{theorem} on separate counters, together with the standard packages those lines need.  The retained environments print in this document as Theorem 1 in order of appearance, whereas the article prints the same items with the numbering of its own full document.  The complete native source of the article as distributed in the arXiv e-print for this exact version is bundled unchanged as \texttt{sources/195164/source0.tex}.
\begin{theorem}
For integers $i,j$, define a sequence of integers $\{ a_{i,j}(n) \}_{n \ge 0}$ recursively by
\begin{align*}
a_{i,j}(n+2) &=
a_{i,j}(n+1) + a_{i-1,j}(n+1) + a_{i,j-1}(n+1) + a_{i-1,j-1}(n+1) \\
&\quad + a_{i-1,j-1}(n),
\end{align*}
with initial conditions
\[
a_{i,j}(0) = 0 \quad \text{for all } i,j,
\]
and
\[
a_{i,j}(1) =
\begin{cases}
1 & \text{if } i=j=1, \\
0 & \text{otherwise}.
\end{cases}
\]
Then
\[
P_{m,m}
=
\left\{ (i,j) \in \mathbb{Z}^2 \,\middle|\,
a_{i,j}(m) \equiv 1 \pmod{2}
\right\}.
\]
\end{theorem}

\begin{proof}
We prove by induction on $m$ that
\[
v_m(i,j) = 0
\quad \Longleftrightarrow \quad
a_{i,j}(m) \equiv 1 \pmod{2}.
\]

For $m=1$, we have $v_1(1,1)=0$, so the claim holds.

Assume that the statement holds for all integers up to $m=n$.
For positive integers $a,b$, the following relations are valid:
\[
a \oplus b =
\begin{cases}
(a-1)\oplus b + 1 & \text{if } u(a) < u(b), \\
a \oplus (b-1) + 1 & \text{if } u(a) > u(b), \\
(a-1)\oplus (b-1) & \text{if } u(a) = u(b).
\end{cases}
\]
In particular, the condition $u(a)=u(b)$ is equivalent to
\[
a \oplus (b-1) = (a-1)\oplus b.
\]

Suppose that
\[
v_{n+1}(i+1,j+1)
= i \oplus j \oplus (n-i) \oplus (n-j) = 0.
\]
Reorder $i,n-i$ as $a,b$ so that $u(a) \le u(b)$, and similarly reorder $j,n-j$ as $c,d$ with $u(c) \le u(d)$.
Then $a+b=c+d=n$ and $a\oplus b=c\oplus d$.

\medskip
\noindent
(1) If \(u(a) < u(b)\) and \(u(c) < u(d)\), then
\[
(a-1)\oplus b \oplus (c-1)\oplus d = 0.
\]
In this case, The following equivalence holds:
\[
a\oplus(b-1)=c\oplus(d-1)
\quad \Longleftrightarrow \quad
(a-1)\oplus(b-1)=(c-1)\oplus(d-1).
\]
Moreover, assuming \((a-1) \oplus b = c \oplus (d-1)\), it follows that
\((c-1) \oplus d = c \oplus (d-1)\), which contradicts \(u(c) < u(d)\).
Therefore,
\[
(a-1) \oplus b \oplus c \oplus (d-1) \neq 0.
\]
Similarly, we have
\[
a \oplus (b-1) \oplus (c-1) \oplus d \neq 0.
\]


\medskip
\noindent
(2) If $u(a)=u(b)$ or $u(c)=u(d)$, then necessarily
\[
u(a)=u(b)=u(c)=u(d),
\]
since the carry positions in the binary addition of $a$ and $b$ coincide with those in the binary addition of $c$ and $d$.
Thus,
\[
(a-1)\oplus(b-1)\oplus(c-1)\oplus(d-1)=0,
\]
and all four of the following expressions vanish:
\[
(a-1)\oplus b \oplus (c-1)\oplus d,
\quad
(a-1)\oplus b \oplus c \oplus (d-1),
\]
\[
a\oplus(b-1)\oplus(c-1)\oplus d,
\quad
a\oplus(b-1)\oplus c \oplus (d-1).
\]

By cases (1), (2), and the induction hypothesis, we conclude that
\[
v_{n+1}(i+1,j+1)=0
\quad \Longleftrightarrow \quad
a_{i+1,j+1}(n+1) \equiv 1 \pmod{2},
\]
which completes the proof.
\end{proof}

\end{document}
